
\subsubsection{6.1 Interpreting the BERT Result}

The $r = 0.984$ correlation exceeds the pre-registered threshold ($r \geq 0.65$) and the pre-experimental confidence estimate (0.72). Three factors contribute:

\textbf{Bimodal oracle structure amplifies measured correlation.} The oracle assigns only $r^* = 4$ or $r^* = 64$ across the 5 measured layers. erank separates these two tiers --- the empirical problem reduces to binary structural classification. Whether erank discriminates intermediate oracle ranks (e.g., $r^* = 8, 16, 32$) cannot be assessed from the current sample.

\textbf{Erank values align with layer-type hierarchy.} FFN intermediate layers (erank $\approx$ 704--710) receive oracle rank 64; attention layers including both output projections and a query projection (erank $\approx$ 557--590) receive oracle rank 4. The structural scalar predicts oracle rank without observing any task data.

\textbf{Sample size caveat.} With $n = 5$ oracle measurements spanning both tiers, the near-perfect Pearson $r$ reflects near-deterministic bimodal discrimination. Full coverage of all 72 layers is necessary to assess whether the correlation persists over intermediate erank values.

\subsubsection{6.2 Global Hypothesis Status}

The pre-registered gate for H-E1 requires $r \geq 0.65$ with $p < 0.05$ in $\geq 2/3$ of the three model families. BERT-base-uncased satisfies this threshold; DeBERTa-v3-base and ViT-base-patch16-224 oracle sweeps have not been completed. The global hypothesis verdict is \textbf{PENDING}. The global \texttt{summary.json} reports \texttt{''global\textbackslash \_pass'': false, ''global\textbackslash \_verdict'': ''FAIL''} for the current run, reflecting that only 1 of 3 required families has been evaluated. These results are presented as preliminary single-family evidence.

\subsubsection{6.3 Limitations}

\textbf{L1: Partial oracle coverage.} 5 of 72 BERT-base-uncased target layers have been oracle-measured (6.9\%). Intermediate erank layers ($\approx 620$--680) are unmeasured; non-linearities may exist.

\textbf{L2: Marginal vs. joint oracle.} The PARA oracle assigns rank per layer independently, holding all other layers at $r = 8$. Globally optimal joint rank allocation --- optimizing all layers simultaneously --- may differ from the marginal oracle assignments.

\textbf{L3: Encoder-only models only.} Decoder-only LLMs (LLaMA, Mistral) with causal attention have different architectural properties; erank-oracle correlation is untested for these architectures.

\textbf{L4: Classification tasks only.} The oracle requires a scalar accuracy metric; extension to generation tasks requires different oracle definitions.

\textbf{L5: erank CV below A1 threshold for NLP encoders.} BERT-base-uncased and DeBERTa-v3-base show CV $\approx$ 0.04, below the pre-registered A1 assumption threshold (CV > 0.05). Only ViT-base-patch16-224 (CV $\approx$ 0.12) meets A1. The strong BERT oracle correlation is driven by oracle bimodality rather than erank spread. If the oracle produces intermediate rank assignments across the full 72 layers, the modest erank variation (range ratio 1.$34\times$--1.$40\times$ for NLP encoders) may limit discriminative power as a continuous predictor.

\textbf{L6: Depth confound in oracle sample.} The 5 oracle-measured BERT layers span depths 0--10, with attention layers sampled from shallower positions (layers 0, 3, 6) and FFN layers from deeper positions (layers 8, 10). Layer type is partially confounded with depth in the current sample. Partial correlation analysis controlling for depth is not feasible at $n = 5$ and is planned once full-coverage oracle sweeps are completed.

\textbf{L7: Stable rank oracle correlation not evaluated.} Stable rank ($\|W_0\|_F^2 / \|W_0\|_2^2$) maps were computed for all BERT layers but oracle correlation for stable rank is not reported, as oracle sweeps covered only 5 layers and this paper focuses on erank as the primary predictor. Stable rank vs. erank oracle comparison is deferred to the full-oracle experiment.

\textbf{L8: Bootstrap CI failure.} Bootstrap confidence interval computation returned NaN at $n = 5$ due to degenerate bimodal resampling. The statistical significance claim rests on the exact one-tailed $p$-value ($p = 0.0013$) from the $t$-distribution.

\subsubsection{6.4 Broader Implications}

If the erank-oracle correlation generalizes across model families and oracle rank distributions, it would enable zero-cost per-layer rank allocation from pretrained weights alone. However, the current evidence --- one model family, 5 oracle measurements, bimodal oracle --- is insufficient to support such a claim. The present findings motivate but do not establish generalization.

The finding is consistent with prior observations that intrinsic dimensionality varies by layer type \citep{Aghajanyan2021Intrinsic} and that structural properties of pretrained weights are informative for adaptation [Meng et al., 2024; Balazy et al., 2024]. The structural predictor approach investigated here is distinct from these works in that it tests oracle rank prediction rather than initialization or adapter design.

